\documentclass[11pt]{article}
\usepackage[a4paper,margin=30mm]{geometry}
\usepackage{amsmath,amssymb,amsthm,mathtools}
\usepackage[hidelinks]{hyperref}
\usepackage{microtype}

\newtheorem{theorem}{Theorem}[section]
\newtheorem{proposition}[theorem]{Proposition}
\newtheorem{lemma}[theorem]{Lemma}
\newtheorem{corollary}[theorem]{Corollary}
\theoremstyle{remark}
\newtheorem{remark}[theorem]{Remark}
\newtheorem{problem}[theorem]{Problem}

\newcommand{\Z}{\mathbb Z}
\newcommand{\R}{\mathbb R}
\newcommand{\vTwo}{v_2}
\newcommand{\spec}{\operatorname{spec}}
\newcommand{\per}{\operatorname{per}}

\title{Arithmetic Periodic Phases in Signed Step Circulants:\\
$2$-Adic Period Compression and Sharp Spectral Gaps}
\author{Working research draft --- author metadata to be confirmed}
\date{9 September 2026}

\begin{document}
\maketitle

\begin{abstract}
We study explicit periodic signings of the step-$(1,s)$ operator and the
continuous Bloch spectral edge of the resulting signed circulants.  The
smallest nontrivial periodic sector already has an exact variational
bifurcation: among all sign words of period dividing two, the alternating
phase is the unique minimizer up to translation; it lies below squared edge
$8$ exactly when $s$ is odd.  Its gap admits a complete one-dimensional
variational description and a high-order expansion beginning with
$\pi^2/s^2$.

For even jumps we uncover a different arithmetic mechanism.  For every even
$L\ge4$ and every odd multiplier $2q+1$, a period-$2L$ two-defect flux phase
for the jump $s=L(2q+1)$ has squared Bloch edge strictly below $8$.  The
$2L$-dimensional fiber folds to an $L$-site two-component chain and then to a
fixed $4\times4$ transfer monodromy.  This gives a complete period-compression
theorem: if $\vTwo(s)=k\ge2$, period $2^{k+1}$ suffices, independently of the
odd part of $s$; the case $\vTwo(s)=1$ is covered by a fixed period-eight
phase with exact edge $4+\sqrt{10+2\sqrt5}$.

The compressed family has a uniform quadratic gap
\[
 8-R_{L,q}\ge \frac1{10L^2},
\]
and this scale is sharp.  Uniformly over the odd multiplier,
\[
 L^2(8-R_{L,q})\longrightarrow
 4\arccos^2\frac13.
\]
The constant arises from the Robin quantization law $3\cos a=1$, rather than
the $\pi/2$ quantization producing the independent $\pi^2$ constant in the
older long-period family.  We also prove first-order phase rigidity and that,
for all sufficiently large even $L$, the unique global maximizing Bloch phase
is exactly $z=1$.  Consequently the compressed gap has a complete algebraic
large-$L$ expansion obtained from a scalar Robin equation.
\end{abstract}

\noindent\textbf{Keywords:} signed graph; circulant graph; periodic operator;
Floquet theory; chiral symmetry; transfer matrix; spectral radius; $2$-adic
valuation

\medskip
\noindent\textbf{MSC 2020:} 05C50 (primary); 05C22, 47B39, 15A18 (secondary)

\section{Introduction and main results}

For a periodic sign word $\tau$ and an integer jump $s\ge2$, consider the
real self-adjoint operator
\begin{equation}
 (A_{s,\tau}x)_j
 =x_{j-1}+x_{j+1}
 +\tau_{j-s}x_{j-s}+\tau_jx_{j+s}.
 \label{eq:operator}
\end{equation}
If $\tau$ has period $p$, let $H_{s,\tau}(z)$ be its $p\times p$ Bloch
fiber under
\[
 x_{j+p}=zx_j,\qquad |z|=1,
\]
and define the continuous squared Bloch edge
\begin{equation}
 R_s(\tau)=\max_{|z|=1}\rho(H_{s,\tau}(z))^2.
 \label{eq:bloch-edge}
\end{equation}
The reference value $8$ is the natural Ramanujan-type squared threshold for
the four-regular support considered here.  Our problem is deliberately
periodic: we ask how efficiently explicit flux phases can force
$R_s(\tau)<8$, how short the period may be, and what asymptotic boundary
conditions govern the resulting gaps.

The answer has two different mechanisms.  Odd jumps are already resolved in
the period-two sector.  Even jumps require defects, but their necessary
period is controlled not by the size of $s$ itself, rather by its $2$-adic
scale.

\subsection{The period-two variational problem}

\begin{theorem}[Exact period-two phase diagram]
\label{thm:period-two}
Among all Hamilton-gauge words whose period divides two, the alternating
words $\tau_j=\pm(-1)^j$ are, up to translation, the unique minimizers of
$R_s(\tau)$.  If $s$ is odd, then
\[
 \min_{\per(\tau)\mid2}R_s(\tau)=8-g_s<8,
 \qquad
 g_s=4\min_{\theta\in\R}
 \bigl(\sin^2\theta+\cos^2(s\theta)\bigr).
\]
If $s$ is even, then
\[
 \min_{\per(\tau)\mid2}R_s(\tau)=8.
\]
For odd $s\ge3$, the minimizing phase $\theta_s\in(0,\pi/(2s))$ is unique
and satisfies
\[
 \sin(2\theta_s)=s\sin(2s\theta_s).
\]
\end{theorem}

\begin{theorem}[Odd-jump expansion]
\label{thm:odd-expansion}
As odd $s\to\infty$,
\[
\begin{aligned}
 g_s={}&\frac{\pi^2}{s^2}
 -\frac{\pi^2(\pi^2+12)}{12s^4}
 +\frac{\pi^2(\pi^4+120\pi^2+360)}{360s^6}\\
 &-\frac{\pi^2(\pi^6+896\pi^4+18480\pi^2+20160)}{20160s^8}
 +O(s^{-10}).
\end{aligned}
\]
Moreover $g_s$ decreases strictly and $s^2g_s$ increases strictly to
$\pi^2$ along odd integers.
\end{theorem}

Thus the first parity obstruction is variational, not merely a feature of a
chosen example.

\subsection{A universal two-defect compression theorem}

Let $L\ge4$ be even.  On residues modulo $2L$, define the local flux word
\begin{equation}
 Q_0=Q_2=1,
 \qquad Q_j=-1\quad(j\ne0,2),
 \label{eq:two-defect-Q}
\end{equation}
and choose the Hamilton-gauge lift with $\tau_0=1$.  Explicitly,
\begin{equation}
 \tau_0=\tau_1=1,
 \qquad \tau_2=\tau_3=-1,
 \qquad \tau_j=(-1)^j\quad(4\le j<2L).
 \label{eq:two-defect-tau}
\end{equation}
For a jump
\[
 s=L(2q+1),\qquad q\ge0,
\]
write $R_{L,q}$ for the squared Bloch edge of this period-$2L$ phase and
$g_{L,q}=8-R_{L,q}$.

\begin{theorem}[General two-defect compression]
\label{thm:compression}
For every even $L\ge4$ and every $q\ge0$,
\[
 \boxed{R_{L,q}<8.}
\]
For every even $L\ge6$,
\[
 \boxed{g_{L,q}\ge\frac1{10L^2}}
\]
uniformly in $q$.
\end{theorem}

The period-$2L$ fiber admits two successive reductions: first to an
$L$-site block Jacobi chain, and then to a fixed $4\times4$ transfer
monodromy.  The proof is therefore uniform in $L$ despite the growing cell.

\begin{corollary}[$2$-adic period compression]
\label{cor:2adic}
Let $s$ be even.  If $\vTwo(s)=1$, a fixed period-eight phase has the exact
edge
\[
 \boxed{4+\sqrt{10+2\sqrt5}<8.}
\]
If $\vTwo(s)=k\ge2$, then an explicit phase of period
\[
 \boxed{2^{k+1}}
\]
has squared Bloch edge below $8$.  Thus the required period depends only on
$\vTwo(s)$ and is independent of the odd part of $s$.
\end{corollary}

This is the main arithmetic compression statement.  For example, a jump
$s=2^k m$ with $m$ large and odd is treated using a cell depending on $2^k$
only, rather than a cell growing linearly with $s$.

\subsection{Sharp compressed gap and phase selection}

Set
\[
 a=\arccos\frac13.
\]

\begin{theorem}[Sharp global compressed constant]
\label{thm:sharp-compressed}
For every sequence of even $L\to\infty$ and every sequence $q=q(L)\ge0$,
\[
 \boxed{
 L^2g_{L,q}\longrightarrow4a^2
 =4\arccos^2\frac13.
 }
\]
Equivalently, the convergence is uniform in the odd multiplier $2q+1$.
If $z_L=e^{it_L}$ is any maximizing Bloch phase and
\[
 d_L=2\cos((2q+1)t_L),
 \qquad e_L=2\cos t_L,
\]
then
\[
 L^2(2-d_L)\to0,
 \qquad e_L\to2.
\]
\end{theorem}

\begin{theorem}[First correction and phase rigidity]
\label{thm:first-correction}
Uniformly in $q$,
\[
 \boxed{
 L^2g_{L,q}
 =4a^2+\frac{16a^2}{3L}+o(L^{-1}).
 }
\]
For every maximizing phase,
\[
 L^3(2-d_L)\to0,
 \qquad
 L(2-e_L)\to0.
\]
\end{theorem}

\begin{theorem}[Eventual exact periodic phase]
\label{thm:eventual-z1}
There exists an even integer $L_0$ such that for every even $L\ge L_0$ and
every $q\ge0$, the unique global maximizing Bloch phase is
\[
 \boxed{z=1.}
\]
Consequently
\[
 R_{L,q}=\rho(H_{L,q}(1))^2
\]
for all such $L$ and $q$.
\end{theorem}

At $z=1$, the fiber is independent of $q$.  If $L=2r$, its top squared
eigenvalue is $6+2\cos\theta_r$, where the unique soft angle satisfies the
scalar Robin equation
\begin{equation}
 3\cos(r\theta_r)
 +2\tan(\theta_r/2)\sin(r\theta_r)=1.
 \label{eq:robin}
\end{equation}
Theorem~\ref{thm:eventual-z1} therefore converts the large-$L$ global Bloch
problem into a one-dimensional analytic equation.

\begin{corollary}[All-order compressed expansion]
\label{cor:all-order}
For every fixed $M\ge0$, there are explicit constants
$\Gamma_0,\ldots,\Gamma_M$ such that, uniformly in $q$,
\[
 L^2g_{L,q}
 =\sum_{j=0}^{M}\Gamma_jL^{-j}+O_M(L^{-M-1})
\]
for all sufficiently large even $L$.  The first four are
\[
 \Gamma_0=4a^2,
 \qquad
 \Gamma_1=\frac{16a^2}{3},
\]
\[
 \Gamma_2
 =\frac{16a^2}{3}
 +\frac{4\sqrt2\,a^3}{9}
 -\frac{4a^4}{3},
\]
and
\[
 \Gamma_3
 =\frac{16a^2}{81}
 \left(24+6\sqrt2\,a-13a^2\right).
\]
\end{corollary}

\subsection{A second spectral constant}

The compressed family is not the only explicit family of interest.  The
older period-$4s$ construction for even $s$ has a different soft boundary
condition and satisfies
\[
 s^2\widehat g_s\to\pi^2.
\]
For $s=2r$, its global maximizing phase has a genuine boundary-layer slip,
\[
 r^2\phi_r\to\frac{\pi}{4\sqrt2},
 \qquad
 r^4(e_r-\widehat g_{2r})\to\frac{\pi^2}{32}.
\]
Thus the constants
\[
 \pi^2
 \qquad\text{and}\qquad
 4\arccos^2\frac13
\]
come from two genuinely different periodic mechanisms.

\section{Switching, flux coordinates, and Bloch fibers}
\label{sec:prelim}

We work in Hamilton gauge, with all nearest-neighbor edges positive and the
remaining chord signs encoded by the word $\tau$.  The local cycle flux is
$Q_j=\tau_j\tau_{j+1}$.  Switching changes vertex signs but leaves the cycle
flux data invariant.  Periodicity of $\tau$ produces the usual Bloch
fibers $H_{s,\tau}(z)$.

The two-defect word \eqref{eq:two-defect-tau} satisfies the reversal identity
\begin{equation}
 \boxed{\tau_{3-j}=-\tau_j}
 \label{eq:reversal}
\end{equation}
modulo $2L$.  Let
\[
 (Dx)_j=(-1)^jx_j,
 \qquad
 (Rx)_j=x_{L+3-j}.
\]
Because $s=L(2q+1)$ and $L$ is even, reflection negates every chord sign,
while $D$ negates every nearest-neighbor term.  Hence
\[
 (DR)A_{s,\tau}(DR)^{-1}=-A_{s,\tau}.
\]
On a single Bloch fiber, reflection reverses the Bloch multiplier; combining
it with complex conjugation gives an antiunitary chiral symmetry.  Therefore
\begin{equation}
 \spec H_{s,\tau}(z)=-\spec H_{s,\tau}(z),
 \label{eq:chiral-spectrum}
\end{equation}
and the characteristic polynomial is a monic polynomial in $\lambda^2$.

\section{The period-two variational problem}
\label{sec:period-two}

For the alternating word and odd $s$, direct squaring cancels the mixed
channels:
\[
 A_s^2=4I+T^2+T^{-2}-T^{2s}-T^{-2s}.
\]
On a Fourier mode $T=e^{i\theta}$,
\[
 F_s(\theta)=4+2\cos(2\theta)-2\cos(2s\theta),
\]
so
\[
 8-\max_\theta F_s(\theta)
 =4\min_\theta
 \bigl(\sin^2\theta+\cos^2(s\theta)\bigr).
\]
The derivative equation gives Theorem~\ref{thm:period-two}.  The strict
uniqueness follows from monotonicity of
$\sin(sx)/\sin x$ on $(0,\pi/s)$.  Analytic inversion of the critical
equation yields Theorem~\ref{thm:odd-expansion}.

The other period-dividing-two words are classified directly.  For even $s$
the alternating word cannot cross below $8$, proving the exact parity
bifurcation in the smallest nontrivial sector.

\section{Folding the two-defect family}
\label{sec:folding}

Choose $\eta$ with $\eta^2=z$ and set
\[
 \omega=z^q\eta=e^{i\beta},
 \qquad c=\cos\beta,
 \qquad r=\sin\beta.
\]
Pair residues $j$ and $j+L$.  After multiplying the second coordinate in
each pair by $\eta^{-1}$ and applying a staggered Pauli gauge, the fiber
becomes an $L$-site block Jacobi matrix.  The interior hopping block is
$\sigma_z$ and the onsite blocks are
\begin{equation}
 V_j=\begin{cases}
 2r\,\sigma_y,&j=1,2,\\
 2c\,\sigma_x,&j\ne1,2.
 \end{cases}
 \label{eq:onsite-blocks}
\end{equation}
The closing block is $i\eta\sigma_y$.

For an eigenvector $\psi_j\in\mathbb C^2$, set
\[
 T(A)=\begin{pmatrix}A&-I_2\\I_2&0\end{pmatrix},
\]
\[
 A_g=\lambda\sigma_z-2ic\sigma_y,
 \qquad
 A_d=\lambda\sigma_z+2ir\sigma_x.
\]
The key scalar-square identities are
\[
 A_g^2=(\lambda^2-4c^2)I_2,
 \qquad
 A_d^2=(\lambda^2-4r^2)I_2.
\]
Hence the full growing-dimensional problem reduces to
\begin{equation}
 M(\lambda)=T(A_g)^{L-3}T(A_d)^2T(A_g),
 \label{eq:monodromy}
\end{equation}
a fixed $4\times4$ monodromy.

\section{Threshold positivity and general period compression}
\label{sec:compression-proof}

Put $y=\lambda^2$.  At $y=8$, define
\[
 d=\omega^2+\omega^{-2}\in[-2,2],
 \qquad
 e=z+z^{-1}\in[-2,2],
\]
\[
 m=\frac{L-4}{2},
 \qquad
 t=\frac{4-d}{2},
 \qquad
 u=U_m(t),
 \qquad
 w=U_{m-1}(t).
\]
The scalar-square relation gives the Chebyshev identity
\[
 u^2+w^2-(4-d)uw=1.
\]
Substitution into \eqref{eq:monodromy} yields
\begin{equation}
 P_{L,q,z}(8)=F_L(d)+d-e,
 \label{eq:threshold-discriminant}
\end{equation}
where
\[
 F_L(x)=u_x^2A(x)+u_xw_xB(x)+C(x),
\]
\[
 A(x)=x^4-21x^2-8x+84,
\]
\[
 B(x)=x^3+4x^2-4x-16,
 \qquad
 C(x)=x^2+11x+6.
\]
The crucial estimate is
\begin{equation}
 \boxed{F_L(x)\ge20\qquad(-2\le x\le2).}
 \label{eq:F-lower}
\end{equation}
It follows from an exact factorization of $F_{m+1}-F_m$ and the base case
$L=4$.  Hence
\[
 P_{L,q,z}(8)\ge16.
\]
At $z=1$ the determinant factors as
\[
 \det(\lambda I-H_{L,q}(1))
 =p_r(y)(p_r(y)+4),
 \qquad r=L/2,
\]
with
\[
\begin{aligned}
 p_r(y)={}&
 U_{r-2}\!\left(\frac{y-6}{2}\right)(y^2-8y+6)\\
 &-U_{r-3}\!\left(\frac{y-6}{2}\right)(y-2)-2.
\end{aligned}
\]
For $y\ge8$,
\[
 p_r(y)+2
 =6(u-w)+(y-8)(yu-w)>0,
\]
which fixes the inertia.  Continuity in $z$ and
\eqref{eq:threshold-discriminant} prove Theorem~\ref{thm:compression}.

To upgrade strict positivity to the scale $L^{-2}$, differentiate the full
characteristic polynomial at $y=8$.  The logarithmic derivative of
$U_m(t)$ is controlled by its zeros in $(-1,1)$, giving
\[
 \frac{U_m'(t)}{U_m(t)}
 \le \frac{m(m+2)}3,
 \qquad t\ge1.
\]
After accounting for $dt/dy=1/2$ and using
\[
 \frac{P_y(8)}{P(8)}
 =\sum_j\frac1{8-y_j},
\]
one obtains
\[
 8-\rho(H_{L,q}(z))^2
 \ge\frac1{10L^2}
\]
uniformly in $z$ and $q$.

\section{Exact Floquet discriminant and the sharp compressed constant}
\label{sec:sharp-compressed}

For general $y$, retain
\[
 d=z^{2q+1}+z^{-(2q+1)},
 \qquad e=z+z^{-1},
\]
\[
 t=\frac{y-d-4}{2},
 \qquad u=U_m(t),
 \qquad w=U_{m-1}(t).
\]
Using the Chebyshev identity to remove $u^2$, the full characteristic
equation takes the exact scalar form
\begin{equation}
 \boxed{\det(\lambda I-H_{L,q}(z))=G_L(y,d)-e.}
 \label{eq:full-discriminant}
\end{equation}
One convenient representation is
\[
 G_L(y,d)=A_*(y,d)uw+B_*(y,d)w^2+C_*(y,d),
\]
where
\[
\begin{aligned}
A_*={}&(-d^2+y^2-8y+10)\\
&\times(d^3-d^2y+4d^2-dy^2+8dy-12d
+y^3-12y^2+40y-32),
\end{aligned}
\]
\[
 B_*=-(-d^2-d+y^2-9y+14)
 (-d^2+d+y^2-7y+6),
\]
and
\[
\begin{aligned}
 C_*={}&d^4-2d^2y^2+16d^2y-20d^2+4d\\
 &+y^4-16y^3+84y^2-160y+90.
\end{aligned}
\]

Let a maximizing root have
\[
 y=8-\frac\gamma{L^2},
 \qquad
 d=2-\frac D{L^2}.
\]
The periodic phase provides $\gamma=O(1)$.  The exact discriminant first
forces $D=O(1)$.  After subsequence extraction, Chebyshev scaling gives
\[
 \frac{u}{L},\frac{w}{L}\to
 \Phi(\gamma-D),
\]
where
\[
 \Phi(\xi)=
 \begin{cases}
 \dfrac{\sin(\sqrt\xi/2)}{\sqrt\xi},&\xi>0,\\[1ex]
 \dfrac12,&\xi=0,\\[1ex]
 \dfrac{\sinh(\sqrt{-\xi}/2)}{\sqrt{-\xi}},&\xi<0.
 \end{cases}
\]
The limiting discriminant is
\begin{equation}
 0=34-e_*-36(\gamma-D)\Phi(\gamma-D)^2.
 \label{eq:continuum-discriminant}
\end{equation}
The hyperbolic case is impossible.  In the oscillatory case,
\[
 36\sin^2\left(\frac{\sqrt{\gamma-D}}2\right)=34-e_*
 \ge32.
\]
Together with the periodic-phase upper bound this forces
\[
 D=0,
 \qquad e_*=2,
 \qquad \gamma=4\arccos^2\frac13,
\]
proving Theorem~\ref{thm:sharp-compressed}.

\section{Robin quantization, rigidity, and exact phase selection}
\label{sec:robin}

At $z=1$, the top root of $p_r(y)=0$ can be written
\[
 y=6+2\cos\theta_r,
\]
where \eqref{eq:robin} has a unique solution in $(0,\pi/(2r))$.  Setting
$a_r=r\theta_r$, one obtains
\[
 a_r\to a=\arccos(1/3).
\]
More precisely,
\[
\begin{aligned}
 a_r={}&a+\frac{a}{3r}
 +\frac{a(\sqrt2 a+8)}{72r^2}\\
 &+\frac{a(10a^2+9\sqrt2 a+24)}{648r^3}
 +O(r^{-4}).
\end{aligned}
\]
This yields Corollary~\ref{cor:all-order} after
\[
 g_{L,q}=4\sin^2\left(\frac{a_r}{2r}\right)
\]
once exact phase selection is known.

For the phase-selection argument, introduce
\[
 h=L^{-1},
 \qquad
 D=L^2(2-d),
 \qquad
 E=2-e,
\]
and write a near-edge root as $y=8-\gamma h^2$.  The exact characteristic
equation extends analytically to $h=0$ with limiting function
\[
 \mathcal F(0,\gamma,D,E)
 =32+E-36\sin^2\left(\frac{\sqrt{\gamma-D}}2\right).
\]
At the base point $\Gamma_0=4a^2$,
\[
 \partial_\gamma\mathcal F=-\frac{2\sqrt2}{a}\ne0.
\]
The implicit-function theorem therefore gives a unique local branch
$\gamma=\Gamma(h,D,E)$, with
\[
 \partial_D\Gamma(0,0,0)=1,
 \qquad
 \partial_E\Gamma(0,0,0)=\frac{a}{2\sqrt2}>0.
\]
Hence nonzero phase displacement strictly increases the gap in a fixed local
neighborhood.  The sharp global theorem puts every maximizing phase in this
neighborhood for sufficiently large $L$.  Comparing with the admissible
periodic phase forces $D=E=0$, proving Theorem~\ref{thm:eventual-z1}.

\section{Exact low layers}
\label{sec:low-layers}

The general compression theorem has especially explicit low-dimensional
models.

\subsection{$\vTwo(s)=1$}
For $s\equiv2\pmod4$, the fixed period-eight word
\[
 (1,1,-1,1,-1,-1,1,-1)
\]
has exact squared Bloch edge
\[
 4+\sqrt{10+2\sqrt5}.
\]
The proof reduces to a $4\times4$ chiral squared block and a square-sum
inequality.

\subsection{$\vTwo(s)=2$}
For $s\equiv4\pmod8$, another fixed period-eight word has top squared branch
\[
 4+\sqrt{
 8+2\cos(2nt)+\sqrt{10-8\cos(nt)+2\cos t}},
 \qquad n=s/4,
\]
and remains strictly below $4+\sqrt{10+2\sqrt5}$.

\subsection{$\vTwo(s)=3$}
For $s\equiv8\pmod{16}$, the period-sixteen specialization has an exact
threshold determinant $F(d)+d-e$ with a positive Bernstein-basis certificate
for $F$ on $[-2,2]$.  It provides a concrete finite-dimensional model of the
general transfer proof.

\section{The second family: the $\pi^2$ law and phase slip}
\label{sec:old-family}

For comparison, retain the older all-jump construction.  Odd $s$ uses the
period-two alternating phase.  Even $s$ uses the primitive period-$4s$
antipodal family.  For this particular parity-dependent family,
\[
 s^2\widehat g_s\to\pi^2.
\]
When $s=2r$, the phase-zero endpoint is not always globally optimal.  The
global phase satisfies
\[
 r^2\phi_r\to\frac{\pi}{4\sqrt2},
\]
and the gain over phase zero is
\[
 r^4(e_r-\widehat g_{2r})\to\frac{\pi^2}{32}.
\]
The compressed family has the opposite phase-selection behavior:
Theorem~\ref{thm:eventual-z1} eventually pins the optimizer exactly at
$z=1$.

\section{Open problems}
\label{sec:open}

The existence hierarchy and the sharp compressed leading asymptotic are now
closed.  The remaining problems concern optimality and finite-size rigidity.

\begin{problem}[Exact finite phase selection]
Prove that $z=1$ is the unique global maximizing Bloch phase of the compressed
two-defect family for every even $L\ge6$, removing the unspecified threshold
$L_0$ from Theorem~\ref{thm:eventual-z1}.
\end{problem}

\begin{problem}[Minimal compressed period]
For $\vTwo(s)=k\ge2$, determine whether period $2^{k+1}$ is minimal inside a
natural two-defect or signed-reflection chiral class.
\end{problem}

\begin{problem}[Fixed-defect variational optimality]
At fixed even $L$, compare the two-defect word with all bounded-defect flux
patterns of period $2L$.  Is the two-defect phase variationally optimal in a
natural defect class?
\end{problem}

\begin{problem}[Sharp finite quadratic constant]
Replace the convenient bound $1/(10L^2)$ by the best finite-$L$ uniform lower
bound and determine its relation to the asymptotic constant
$4\arccos^2(1/3)$.
\end{problem}

\end{document}
